\section{Conclusions}

\begin{frame}{Conclusions}
	In this course, we reviewed probability fundamentals:
	\begin{itemize}
		\item Conditional probability.
		\item Bayes' theorem.
		\item Bayes factor.
		\item Recursive nature of Bayes' theorem.
	\end{itemize}
	\pause
	We also saw how the Naive Bayes algorithm utilizes Bayes' theorem to solve classification problems:
	\begin{itemize}
		\item With data following a multinomial distribution (text data): Multinomial Naive Bayes algorithm.
		\item With data following a Gaussian distribution: Gaussian Naive Bayes algorithm.
	\end{itemize}
	\pause
	We also saw that the naive conditional independence assumption allows us to handle classification problems with input features following different types of distributions (multinomial, Gaussian, etc.) by combining multiple Naive Bayes algorithms.
\end{frame}

\begin{frame}{Conclusions}
	Finally, we introduced rudimentary text data processing techniques:
	\begin{itemize}
		\item Tokenization and stop words.
		\item Lemmatization and Part-of-Speech (POS) tagging.
		\item Count vectorization and TF-IDF.
	\end{itemize}
	\pause
	These processing steps are not exclusive to Naive Bayes algorithms and can be reused for other machine learning algorithms.
\end{frame}